
% 第 XII 篇，第 4--6 节。
\section{约化中心}
\SGARefTarget{sga3:xii:section:4}
\paragraph{\textbf{定义 4.1。}}
\SGARefTarget{sga3:xii:4.1}
设 \(S\) 是预概形，\(G\) 是在 \(S\) 上有限呈示且纤维仿射的 \(S\)-群，\(Z\) 是 \(G\) 的群子预概形。若满足以下条件，则称 \(Z\) 是 \(G\) 的\emph{约化中心}：
\begin{enumerate}
\item[(i)] \(Z\) 是中心的且为乘法型；
\item[(ii)] 对每个基变换 \(S'\to S\) 以及每个中心同态
\[
  u:H\longrightarrow G_{S'},
\]
其中 \(H\) 是 \(S'\) 上的有限型乘法型群，同态 \(u\) 都经由 \(Z_{S'}\) 分解。
\end{enumerate}
（不假设 \(G\) 的纤维仿射时，这个概念的一个变体见 \SGARefLink{sga3:xii:8.6}{8.6}。）
这样的 \(Z\) 必在 \(S\) 上有限型，因为其纤维有限型，故它作为 \(G\) 中的最大中心乘法型子群而唯一确定。容易给出这样的例子：\(G\) 在 \(S\) 上平滑且仿射，存在最大的中心乘法型子群 \(Z\)，但 \(Z\) 不是约化中心，即 \(G\) 不存在约化中心；例如见 \SGARefLink{sga3:xii:1.6}{1.6}。不过，由 \SGARefLink{sga3:ix:theorem:6:8}{IX~6.8} 容易得到，\(G\) 的子群 \(Z\) 是约化中心，当且仅当它是最大的中心乘法型子群，且在每次基变换后仍保持此性质。另见下文 \SGARefLink{sga3:xii:4.3}{4.3}。
由 \SGARefLink{sga3:xii:4.1}{4.1} 显然可知，若 \(Z\) 是 \(G\) 的约化中心，则在每次基变换 \(S'\to S\) 后，\(Z_{S'}\) 是 \(G_{S'}\) 的约化中心。结合此事实、约化中心的唯一性及忠实平坦拟紧下降（SGA~1，VIII），得到：
\paragraph{\textbf{命题 4.2。}}
\SGARefTarget{sga3:xii:4.2}
\emph{设 \(G\) 是在 \(S\) 上有限呈示且纤维仿射的 \(S\)-群预概形。若 \(G\) 对忠实平坦拟紧拓扑局部地存在约化中心，则它存在约化中心。\(Z\) 是 \(G\) 的约化中心的充要条件是，它对 fpqc 拓扑局部地如此。}
还记下：
\paragraph{\textbf{命题 4.3。}}
\SGARefTarget{sga3:xii:4.3}
\emph{设 \(G\) 是有限呈示且纤维仿射的 \(S\)-群预概形，\(Z\) 是群子预概形。\(Z\) 是 \(G\) 的约化中心的充要条件是，它为乘法型，且对每个 \(s\in S\)，\(Z_s\) 是 \(G_s\) 的约化中心。}
事实上，\SGARefLink{sga3:ix:theorem:5:6:item:b}{IX~5.6(b)} 首先表明 \(Z\) 是中心的。由于所考虑的条件在基变换下稳定，只须再证明每个中心同态
\[
  u:H\longrightarrow G,
\]
（其中 \(H\) 为有限型乘法型群）都经由 \(Z\) 分解。由于 \(Z\) 是中心的，\(u\) 与典范浸入 \(v:Z\to G\) 定义群同态
\[
  w:H\times_S Z\longrightarrow G.
\]
由 \SGARefLink{sga3:ix:theorem:6:8}{IX~6.8}，此同态有群像 \(K\)，它是 \(G\) 的乘法型子群。一切归结为证明 \(K=Z\)。由关于 \(Z\) 的假设，逐纤维地这已成立，现在应用 \SGARefLink{sga3:ix:theorem:5:1:bis}{IX~5.1 bis}，即完成 \SGARefLink{sga3:xii:4.3}{4.3} 的证明。
在判据 \SGARefLink{sga3:xii:4.3}{4.3} 中，对每个 \(s\in S\)，\(Z_s\) 是 \(G_s\) 的约化中心这一假设，实际上是纯几何的：只须在取 \(\kappa(s)\) 的代数闭包后验证它，这来自 \SGARefLink{sga3:xii:4.2}{4.2} 的第二个断言。
\paragraph{\textbf{定理 4.4。}}
\SGARefTarget{sga3:xii:4.4}
\emph{设 \(G\) 是域 \(k\) 上的仿射代数群。则 \(G\) 存在约化中心。}
由于 \(G\) 的中心可由 \(G\) 的闭子群表示（\SGARefLink{sga3:viii:corollary:6:7}{VIII~6.7}），立即可归约到 \(G\) 交换的情形。此外，由 \SGARefLink{sga3:xii:4.2}{4.2}，可设 \(k\) 代数闭。在第 XVII 篇 中将看到，此时 \(G\) 可写为积 \(Z\times U\)，其中 \(Z\) 为乘法型，\(U\)“幺幂”，立即可知 \(Z\) 是 \(G\) 的约化中心。这里给出一个不依赖于刚才所述交换代数群一般结构定理的证明。
由 \SGARefLink{sga3:ix:theorem:6:8}{IX~6.8}，\(G\) 的乘法型子群 \(H\) 所成集合向上滤过。我们证明它有最大元素。
当 \(G\) 在 \(k\) 上平滑时，结构定理 BIBLE~4，定理 4 给出
\[
  G=G_s\times G_u,
\]
其中 \(G_s\) 为乘法型，\(G_u\)“幺幂”：这里是指 \(G_u\) 有一个合成列，其因子为 \(\mathbb G_a\) 的子群，并可要求它们平滑。事实上，由 BIBLE~4，定理 3 的推论，\(G_u/G_u^0\) 是幺幂群，故由 BIBLE~4，命题 4，它是 \(p\)-群，其中 \(p\) 为特征指数。另一方面，由 BIBLE~6，定理 1、推论 1，\(G_u^0\) 有一个相继商为平滑连通一维群的合成列，而由 BIBLE~7，定理 4，这些商同构于 \(\mathbb G_a\)。从乘法型群到 \(\mathbb G_a\)，从而到 \(G_u\) 的每个同态都平凡，这来自下面的引理。因此，\(G\) 的每个乘法型子群都包含于 \(G_s\)。
在一般情形，考虑 \(G\) 的子群
\[
  G_0=G_{\mathrm{red}}
\]
它在 \(k\) 上平滑，故由前一种情形有最大的乘法型子群 \(Z_0\)。\(G\) 中包含 \(Z_0\) 的乘法型子群，对应于 \(G'=G/Z_0\) 的乘法型子群。有正合列
\[
  0\longrightarrow G_0/Z_0
   \longrightarrow G/Z_0
   \longrightarrow G/G_0
   \longrightarrow 0,
\]
其中 \(G_0/Z_0\) 除单位群外不含乘法型子群。由于乘法型群的每个子群仍为乘法型（IX~8），\(G/Z_0\) 的每个乘法型子群 \(H\) 满足
\[
  H\cap(G_0/Z_0)=0;
\]
故 \(H\to G/G_0\) 单射。代数群 \(G/G_0\) 是纯不可分的（radicial）的，因而在 \(k\) 上有限，所以 \(H\) 本身纯不可分的（radicial），且其秩以 \(G/G_0\) 的秩为上界。因此，\(G\) 的乘法型子群族有一个最大元素，记为 \(Z\)。
我们断言，\(G/Z\) 除单位子群外不含乘法型子群。事实上，两个乘法型代数群的交换扩张仍为乘法型（\SGARefLink{sga3:xvii:proposition:7:1:1}{7.1.1}）。若 \(Z'\) 是 \(G/Z\) 中某个乘法型子群在 \(G\) 中的逆像，则 \(Z'\) 为乘法型。\(Z\) 的极大性给出 \(Z'=Z\)，因而 \(Z'/Z=0\)。
现在由“有限扩张原理”容易看出，对 \(k\) 的每个扩张 \(K\)，
\[
  (G/Z)_K=G_K/Z_K
\]
同样除单位群外不含乘法型子群。
现在可以证明 \(Z\) 是 \(G\) 的约化中心。设
\[
  u:H\longrightarrow G_S
\]
是 \(S\)-群同态，其中 \(S\) 是 \(k\) 上的预概形，\(H\) 是 \(S\) 上有限型乘法型群。必须证明 \(u\) 经由 \(Z_S\) 分解，或等价地，复合
\[
  H\longrightarrow G_S\longrightarrow
  (G/Z)_S=G_S/Z_S
\]
平凡。置 \(U=G/Z\)。每个同态 \(H\to U_S\) 都平凡：由 \SGARefLink{sga3:ix:corollary:5:2}{IX~5.2}，可归约到 \(S\) 是域 \(K\) 的谱的情形，再由 \SGARefLink{sga3:ix:theorem:6:8}{IX~6.8}，将断言归结为 \(U_K\) 不含非平凡乘法型子群这一事实。只要证明以下引理，\SGARefLink{sga3:xii:4.4}{4.4} 即得证。
\paragraph{\textbf{引理 4.4.1。}}
\SGARefTarget{sga3:xii:4.4.1}
\emph{设 \(H\) 是乘法型 \(S\)-群预概形。则每个 \(S\)-群同态
\[
  u:H\longrightarrow\mathbb G_a
\]
都平凡。}
将 \(E=\mathcal O_S^2\) 视为一个扩张，其子模为 \(E'\simeq\mathcal O_S\)，商为 \(E''\simeq\mathcal O_S\)。则 \(\mathbb G_a\) 可与此扩张的自同构预概形等同。因此，同态 \(u:H\to\mathbb G_a\) 等价于 \(E\) 上一个保持此扩张的 \(H\)-模结构，其中 \(E'\) 在 \(H\) 下稳定，且在 \(E'\) 和 \(E''\) 上诱导的作用平凡。由 \SGARefLink{sga3:i:proposition:4-7-3}{I~4.7.3}，\(H\) 在 \(E\) 上平凡作用；故 \(u\) 平凡。
\paragraph{\textbf{注 4.4.2。}}
\SGARefTarget{sga3:xii:4.4.2}
若 \(G\) 是 \(k\) 上的非零阿贝尔簇，则去掉“\(G\) 仿射”的限制后，\(G\) 不存在 \SGARefLink{sga3:xii:4.1}{4.1} 意义下的约化中心。事实上，若 \(n\) 与特征互素，则 \({}_nG\) 在 \(k\) 上平展，阶与特征互素，因而为乘法型。各 \({}_nG\) 所成族在 \(G\) 中概形稠密。因此，若存在约化中心，它必等于 \(G\)，但这不可能，因为 \(G\) 不是乘法型。这解释了 \SGARefLink{sga3:xii:4.1}{4.1} 中要求 \(G\) 的纤维仿射的原因。
\paragraph{\textbf{引理 4.5。}}
\SGARefTarget{sga3:xii:4.5}
\emph{设 \(S\) 是预概形，\(G\) 是纤维连通的平滑仿射 \(S\)-群预概形，\(T\) 是 \(G\) 的极大环面，\(u:H\to G\) 是中心同态，其中 \(H\) 是有限型乘法型 \(S\)-群预概形。则 \(u\) 经由 \(T\) 分解。}
令 \(C=\operatorname{Centr}_G(T)\)。由 \SGARefLink{sga3:xi:corollary:5:3}{XI~5.3}，它是 \(G\) 的闭平滑群子预概形，因而在 \(S\) 上仿射。由于 \(T\) 在 \(C\) 中为中心子群，它是不变的，由 \SGARefLink{sga3:viii:theorem:5:1}{VIII~5.1}，商群
\[
  U=C/T
\]
可表。由于 \(u\) 是中心的，它经由 \(C\) 分解，只须证明复合
\[
  H\longrightarrow C\longrightarrow U=C/T
\]
平凡。由 \SGARefLink{sga3:ix:corollary:5:2}{IX~5.2}，归约到 \(S\) 是域的谱的情形，并可设该域代数闭。于是 BIBLE~6，定理 2 表明，\(U\) 是在 \(k\) 上平滑的连通“幺幂”代数群；如前所述，它有一个相继商同构于 \(\mathbb G_a\) 的合成列。因此，从有限型乘法型群 \(H\) 到 \(U\) 的每个同态都平凡。\SGARefLink{sga3:xii:4.5}{4.5} 证毕。
\paragraph{\textbf{推论 4.6。}}
\SGARefTarget{sga3:xii:4.6}
\emph{设 \(G\) 如 \SGARefLink{sga3:xii:4.5}{4.5}。若 \(G\) 存在约化中心，则该中心包含于 \(G\) 的每个极大环面。}
\paragraph{\textbf{定理 4.7。}}
\SGARefTarget{sga3:xii:4.7}
\emph{设 \(S\) 是预概形，\(G\) 是在 \(S\) 上平滑且仿射、纤维连通的 \(S\)-群预概形。}
\begin{enumerate}
\item[\textup{\SGAInlineRefTarget{sga3:xii:4.7:item:a}(a)}] \emph{对每个 \(s\in S\)，用 \(z(s)\) 表示 \SGARefLink{sga3:xii:4.4}{4.4} 所定义的 \(G_s\) 的约化中心的型。在有限型 \(\mathbb Z\)-模的同构类集合 \(E\) 上规定如下序：若 \(M'\) 同构于 \(M\) 的商，则称 \(M\) 的类大于 \(M'\) 的类。于是映射
\[
  S\longrightarrow E,\qquad s\longmapsto z(s),
\]
下半连续。}
\item[\textup{\SGAInlineRefTarget{sga3:xii:4.7:item:b}\SGAInlineRefTarget{sga3:xii:theorem:4-7:item:b}(b)}] \emph{群 \(G\) 存在约化中心 \(Z\)，当且仅当上述函数 \(z:S\to E\) 局部常值。若此条件成立，则 \(G/Z\) 可表\textup{（参见\ \SGARefLink{sga3:viii:theorem:5:1}{VIII~5.1}）}，且其约化中心为单位子群。}
\item[\textup{\SGAInlineRefTarget{sga3:xii:4.7:item:c}\SGAInlineRefTarget{sga3:xii:theorem:4-7:item:c}(c)}] \emph{设 \(G\) 的约化秩局部常值\textup{（参见\ \SGARefLink{sga3:xii:1.7:item:b}{1.7(b)}）}。则 \(G\) 存在约化中心 \(Z\)。若 \(G/Z\) 可表，例如 \(G\) 在 \(S\) 上仿射，则 \(G\) 的极大环面 \(T\)\textup{（分别为 \(G\) 的 Cartan 子群 \(C\)）}与 \(G'=G/Z\) 的极大环面 \(T'\)\textup{（分别为 \(G'\) 的 Cartan 子群 \(C'\)）}之间有一一对应：
\[
  T'=T/Z,\qquad C'=C/Z,
\]
反过来，
\[
  T=\varphi^{-1}(T'),\qquad C=\varphi^{-1}(C'),
\]
其中 \(\varphi:G\to G'\) 是典范同态。}
\item[\textup{\SGAInlineRefTarget{sga3:xii:4.7:item:d}(d)}] \emph{设 \(T\) 是 \(G\) 的极大环面，\(\mathfrak g\) 表示 \(G\) 的 Lie 代数，考虑由 \(G\) 的伴随表示\textup{（II~4）}诱导的同态
\[
  \theta:T\longrightarrow\operatorname{GL}(\mathfrak g)
\]
则 \(\ker\theta\) 是 \(G\) 的约化中心。}
\end{enumerate}
\emph{证明。}
\textup{\SGARefLink{sga3:xii:4.7:item:a}{(a)}} 和\textup{\SGARefLink{sga3:xii:4.7:item:b}{(b)}}。\textup{\SGARefLink{sga3:xii:4.7:item:a}{(a)}} 以及\textup{\SGARefLink{sga3:xii:4.7:item:b}{(b)}} 的第一个断言的证明，本质上与 \SGARefLink{sga3:xii:1.7:item:a}{1.7(a)} 和 \SGARefLink{sga3:xii:4.7:item:b}{(b)} 的证明相同，这里不再重复。只指出，在\textup{\SGARefLink{sga3:xii:4.7:item:a}{(a)}} 中必须使用 \SGARefLink{sga3:ix:theorem:5:6}{IX~5.6}，并用到 \(G\) 的纤维连通这一事实。
证明\textup{\SGARefLink{sga3:ix:theorem:5:6:item:b}{(b)}} 的第二个断言。设 \(Z\) 是 \(G\) 的约化中心，置 \(G'=G/Z\)。由 \SGARefLink{sga3:viii:theorem:5:1}{VIII~5.1}，此商可表。它在 \(S\) 上仿射，由 \SGARefLink{sga3:viii:corollary:5:8}{VIII~5.8} 在 \(S\) 上有限呈示，甚至在 \(S\) 上平滑：由于 \(G'\) 在 \(S\) 上平坦且有限呈示，可逐纤维检验这一点，而 VI\textsubscript{B}~\SGARefLink{concl:sga3-VIB-9.2-xii}{9.2(xii)} 断言，域上的平滑代数群的商平滑。由于 \(G\) 的纤维连通，\(G'\) 亦然。因此，\(G'\) 满足与 \(G\) 相同的假设。
由 \SGARefLink{sga3:xii:4.3}{4.3}，可归约到 \(S\) 是域的谱的情形。设 \(Z'\) 是 \(G'\) 的中心乘法型子群；只须证明 \(Z'\) 为单位子群。令 \(Z_1\) 为 \(Z'\) 在 \(G\) 中的逆像，考虑 \(Z_1\) 通过内自同构在 \(G\) 上的作用。由于 \(Z\) 是中心的，它平凡作用，所以 \(Z_1\) 的作用经由 \(Z'\) 在 \(G\) 上的作用分解。此外，\(Z'\) 在 \(Z\) 及 \(G/Z=G'\) 上的作用都平凡。由于 \(Z\) 是中心的，\(Z'\) 在 \(G\) 上的作用因而来自群同态
\[
  u:Z'\longrightarrow
    \operatorname{Hom}_{S\text{-}\mathrm{gr}}(G',Z),
\]
其刻画为
\[
  r(z')\cdot g
    =g\cdot u(z')\bigl(\varphi(g)\bigr),
\]
其中 \(r:Z'\to\operatorname{Aut}_{S\text{-}\mathrm{gr}}(G)\) 是刚才考虑的表示，\(\varphi:G\to G'\) 是典范同态。给定 \(u\) 等价于给定群同态
\[
  v:G'\longrightarrow
    \operatorname{Hom}_{S\text{-}\mathrm{gr}}(Z',Z).
\]
由 \SGARefLink{sga3:x:corollary:5-8}{X~5.8}，右边的群可表且在 \(S\) 上平展。因此其单位截面是开浸入，所以 \(\ker v\) 是 \(G'\) 的开子群。由于 \(G'\) 的纤维连通，\(\ker v=G'\)。因此 \(v=0\)，故 \(u=0\)，于是 \(Z'\)，从而 \(Z_1\)，都在 \(G\) 上平凡作用。因此 \(Z_1\) 在 \(G\) 中是中心的。它是乘法型群 \(Z'\) 按乘法型群 \(Z\) 的交换扩张，故在域上由 \SGARefLink{sga3:xvii:proposition:7:1:1}{7.1.1} 知其本身为乘法型。由于它在 \(G\) 中是中心的，必包含于约化中心 \(Z\)。于是 \(Z_1=Z\)，故 \(Z'\) 是单位子群。这证明了\textup{\SGARefLink{sga3:xii:4.7:item:b}{(b)}}。
\textup{\SGARefLink{sga3:xii:4.7:item:d}{(d)}} 置 \(Z=\ker\theta\)，这是 \(T\) 的乘法型子群，例如由 \SGARefLink{sga3:ix:theorem:6:8}{IX~6.8} 可知。由 \SGARefLink{sga3:xii:4.5}{4.5}，每个中心同态 \(u:H\to G\)，其中 \(H\) 为有限型乘法型群，都经由 \(T\) 分解，因而经由 \(Z\) 分解。由于 \(Z\) 的构造与每个基变换交换，只须再证 \(Z\) 是中心的，或等价地，其中心化子 \(C\) 等于 \(G\)。由 \SGARefLink{sga3:xi:corollary:5:3}{XI~5.3}，\(C\) 是 \(G\) 的闭平滑子群。由于 \(Z\) 在 \(\mathfrak g\) 上平凡作用，有
\[
  \operatorname{Lie}(C)=\mathfrak g.
\]
由此可得 \(C=G\)。事实上，浸入 \(C\to G\) 平展，这可逐纤维检验：它是同维数平滑代数群之间的非分歧同态\textup{（VI\textsubscript{B}~\SGARefLink{sga3:xii:1.3}{1.3}）}，所以 \SGARefLink{sga3:x:lemma:3:5}{X~3.5} 适用。故 \(C\to G\) 是平展闭浸入，因而由 SGA~1，\SGARefLink{sga3:i:subsection:5-1}{I~5.1}，它是开浸入。由于它既开又闭，且 \(G\) 的纤维连通，它是同构。
\textup{\SGARefLink{sga3:xii:4.7:item:c}{(c)}} 由 1.7(b)，\(G\) 对 fpqc 拓扑局部地存在极大环面。因此，\SGARefLink{sga3:xii:4.2}{命题 4.2} 和第\textup{\SGARefLink{sga3:xii:corollary:6-7:equation:d}{(d)}} 项表明 \(G\) 存在约化中心 \(Z\)。第\textup{\SGARefLink{sga3:xii:corollary:6-7:equation:d}{(d)}} 项还表明，\(G\) 的每个极大环面 \(T\) 都包含 \(Z\)。显然 \(T'=T/Z\) 是环面。为证明它是 \(G'\) 的极大环面，\SGARefLink{sga3:xii:1.3}{定义 1.3} 将我们归约到 \(S\) 是代数闭域的谱的情形；此时断言来自\emph{BIBLE}~7，定理 3\SGARefLink{sga3:xii:4.7:item:a}{(a)}。
反之，设 \(T'\) 是 \(G'\) 的极大环面，置 \(T=\varphi^{-1}(T')\)。证明 \(T\) 是 \(G\) 的极大环面。此问题对 fpqc 拓扑是局部的，所以可设 \(G\) 已有极大环面 \(T_0\)。由前文，\(T'_0=T_0/Z\) 是 \(G'\) 的极大环面。由共轭定理 1.7(c)，\(T'\) 与 \(T'_0\) 对 fpqc 拓扑局部共轭。由于 \(\varphi:G\to G'\) 是此拓扑的覆盖，\(G'\) 的截面局部可提升为 \(G\) 的截面；因此 \(T\) 与 \(T_0\) 也局部共轭。故 \(T\) 是极大环面。对 Cartan 子群的论证相同。证明完毕。
当 \(T\) 可对角化，例如 \(T=D_S(M)\)，其中 \(M\) 是有限型自由 \(\mathbb Z\)-模时，\textup{\SGARefLink{sga3:xii:corollary:6-7:equation:d}{(d)}} 有一个有用的重述。由 \SGARefLink{sga3:i:proposition:4-7-3}{I~4.7.3}，\(T\)-模 \(\mathfrak g\) 分解为其 \(T\)-子模 \(\mathfrak g_m\)（\(m\in M\)）的直和：
\[
  \mathfrak g=\bigoplus_{m\in M}\mathfrak g_m.
\]
这些子模必局部自由。假设对每个 \(m\in M\)，\(\mathfrak g_m\) 的秩常值，特别地当 \(S\) 连通时便如此。则“根”的集合
\[
  R=\{m\in M\mid\mathfrak g_m\ne0\}
\]
有限。得到：
\paragraph{\textbf{推论 4.8。}}
\SGARefTarget{sga3:xii:4.8}
\emph{在上述假设与记号下，\(G\) 的约化中心是 \(T\) 中根特征标 \(m\in R\) 的核的交。因此有同构
\[
  Z\simeq D_S(N),
\]
其中 \(N\) 是 \(M\) 对 \(R\) 生成的子群的商。}
\paragraph{\textbf{推论 4.9。}}
\SGARefTarget{sga3:xii:4.9}
\emph{设 \(G\) 是代数闭域 \(k\) 上的代数群。假设 \(G\) 在 \(k\) 上平滑、连通且仿射，其约化中心为单位子群，且其 Lie 代数 \(\mathfrak g\) 幂零。则 \(G\)“幺幂”，即它有一个相继商同构于 \(\mathbb G_a\) 的合成列。}
由\emph{BIBLE}~6，定理 4、推论 3，只须证明 \(G\) 的极大环面 \(T\) 是单位子群，或等价地，其 Lie 代数 \(\mathfrak t\) 为零。由 \SGARefLink{sga3:i:proposition:4-7-3}{I~4.7.3}，\(T\)-模 \(\mathfrak g\) 按 \(T\) 的特征标分解，所以对每个 \(t\in\mathfrak t\)，\(\mathfrak g\) 的自同态 \(\operatorname{ad}_{\mathfrak g}(t)\) 半单。若 \(\mathfrak g\) 幂零，这个自同态也幂零，故为零。由 \SGARefLink{sga3:xii:4.7:item:d}{4.7(d)}，约化中心的假设使同态
\[
  T\longrightarrow\operatorname{GL}(\mathfrak g)
\]
为单态射，因而在 Lie 代数上诱导单射。由 \SGARefLink{sga3:ii:proposition:4-5}{II~4.5}，这意味着 \(\operatorname{ad}_{\mathfrak g}(t)=0\) 蕴含 \(t=0\)。因此 \(\mathfrak t=0\)，正如所需。
\paragraph{\textbf{命题 4.10。}}
\SGARefTarget{sga3:xii:4.10}
\emph{设 \(G\) 是代数闭域 \(k\) 上的平滑连通仿射代数群。则 \(G\) 的约化中心是 \(G\) 的所有极大环面的交。}
这里的交取概形意义，或等价地取 \(G\) 的子函子的意义：它是包含于每个极大环面的 \(G\) 的最大闭子预概形。由于 \(G\) 是 Noether 的，它已是适当有限族极大环面的交。
用 \(Z\) 表示这个交。它是环面的闭子群，故为乘法型，且由 \SGARefLink{sga3:xii:4.5}{4.5} 包含 \(G\) 的约化中心。为证明等号，只须再证 \(Z\) 是中心的。由于 \(G\) 连通，\SGARefLink{sga3:ix:corollary:5:5}{IX~5.5} 将此归结为证明 \(Z\) 不变。由构造，\(Z\) 在每个内自同构 \(\operatorname{int}(g)\)（\(g\in G(k)\)）下不变。因此，其正规化子 \(N\) 由 \SGARefLink{sga3:viii:corollary:6:7}{VIII~6.7} 可表，是 \(G\) 的闭子群，包含 \(G\) 的所有有理点。由于 \(G\) 约化，\(N=G\)，结果得证。
\paragraph{\textbf{命题 4.11。}}
\SGARefTarget{sga3:xii:4.11}
\emph{设 \(S\) 是预概形，\(G\) 是纤维连通且幺幂秩为零的平滑仿射 \(S\)-群预概形。则由 \SGARefLink{sga3:xii:1.9}{1.9}，\(G\) 具有局部常值约化秩，故 \SGARefLink{sga3:xii:1.10}{1.10} 定义了“\(G\) 的极大环面预概形”，记为 \(\mathcal T\)，它在 \(S\) 上平滑、分离且有限型。令 \(G\) 通过内自同构作用于 \(\mathcal T\)，给出群函子同态
\[
  u:G\longrightarrow\operatorname{Aut}_S(\mathcal T).
\]
在这些假设下，\(G\) 的以下三个子函子相同：
\begin{enumerate}
\item[\textup{(i)}] \(G\) 的约化中心 \(Z\)；
\item[\textup{(ii)}] \(G\) 的中心 \(Z'\)；
\item[\textup{(iii)}] 上述同态的核 \(Z''=\ker u\)。
\end{enumerate}
特别地，\(G\) 的中心可由 \(G\) 的乘法型子群表示。}
\emph{证明。}
显然 \(Z\subset Z'\subset Z''\)。只须再证 \(Z''\subset Z\)，或者，由于假设在基变换下稳定，只须证明 \(G\) 在 \(S\) 上每个在 \(\mathcal T\) 上平凡作用的截面 \(g\)，都是 \(Z\) 的截面。置 \(G'=G/Z\)。由 \SGARefLink{sga3:xii:4.7:item:b}{4.7(b)} 和 \SGARefLink{sga3:xii:4.7:item:c}{(c)}，\(G'\) 的极大环面预概形 \(\mathcal T'\) 典范同构于 \(\mathcal T\)。因此归约到 \(G=G'\) 的情形，此时约化中心 \(Z\) 为单位子群。注意，\SGARefLink{sga3:xii:4.7:item:c}{4.7(c)} 还表明 \(G'\) 的幺幂秩等于 \(G\) 的幺幂秩，故为零。
必须证明 \(g\) 是 \(G\) 的单位截面。通常的归约过程首先归约到 \(S\) 为 Noether，再归约到 \(S\) 为局部 Artin：当 \(S\) 局部 Noether 时，只须在每个 \(S'\) 为局部 Artin 的基变换 \(S'\to S\) 后检验断言。此时，由 \SGARefLink{sga3:viii:proposition:6:2:item:a:01}{VIII~6.2(a)} 和 \SGARefLink{sga3:viii:examples:6:5:item:c:01}{6.5(c)}，\(Z''=\ker u\) 可表；也可由 \SGARefLink{sga3:xi:theorem:6:8}{XI~6.8} 得到可表性。由 Nakayama 引理，只须证明纤维 \(Z''_0\) 为单位子群。因此可设 \(S\) 是域 \(k\) 的谱，并可取 \(k\) 代数闭。
子群 \(Z''\) 位于 \(\mathcal T(k)\) 的每个点的稳定子中，即位于 \(G\) 的每个极大环面 \(T\) 的正规化子 \(N(T)\) 中。由于 \(G\) 的幺幂秩为零，\SGARefLink{sga3:xi:proposition:5:9}{XI~5.9} 表明 \(T\) 是 \(N(T)\) 的开子群。因此
\[
  \operatorname{Lie}N(T)=\operatorname{Lie}T,
\]
且 \(\operatorname{Lie}Z''\) 包含于每个极大环面 \(T\) 的 \(\operatorname{Lie}T\)。\SGARefLink{sga3:xii:4.10}{命题 4.10} 表明，所有极大环面的 Lie 代数之交，是约化中心 \(Z\) 的 Lie 代数，在当前归约情形中为零。故 \(\operatorname{Lie}Z''=0\)，所以 \(Z''\) 在 \(k\) 上平展。
此外，\(Z''\) 在 \(G\) 中不变。由于 \(G\) 连通，\(Z''\) 在 \(G\) 中是中心的。因此，对每个极大环面 \(T\)，它包含于
\[
  T=\operatorname{Centr}_G(T)
\]
从而包含于这些环面的交。由 \SGARefLink{sga3:xii:4.10}{4.10}，此交为 \(Z=0\)，证明完毕。
\section[在乘法型子群预概形上的应用]
{在乘法型子群预概形上的应用\texorpdfstring{\protect\footnotemark}{}}
\footnotetext{本节结果在不假设 \(G\) 在 \(S\) 上仿射的情形中的推广，见 M.~Raynaud，《群概形与齐性空间上的丰沛层》〔“Ample sheaves on group schemes and homogeneous spaces”〕，第 IX 章。}
\SGARefTarget{sga3:xii:5}
\paragraph{\textbf{定理 5.1。}}
\SGARefTarget{sga3:xii:5.1}
\emph{设 \(S\) 是预概形，\(G\) 是平滑仿射 \(S\)-群预概形，\(M\) 是表示 \SGARefLink{sga3:xi:theorem:4:1}{XI~4.1} 所述函子的“\(G\) 的乘法型子群预概形”。对每个整数 \(n>0\)，令 \(T_n\) 为 \(M\) 的子函子，使 \(T_n(S')\) 是 \(G_{S'}\) 中满足下式的乘法型子群预概形 \(H\) 的集合
\[
  n\operatorname{id}_H=0.
\]
于是由 \SGARefLink{sga3:xi:proposition:3:12:item:a}{XI~3.12(a)}，\(T_n\) 可表且在 \(S\) 上仿射。令
\[
  u_n:M\longrightarrow T_n
\]
为由 \(u_n(H)={}_nH\) 定义的态射，其中 \({}_nH=\ker(n\operatorname{id}_H)\)。则：}
\begin{enumerate}
\item[\textup{\SGAInlineRefTarget{sga3:xii:theorem:5-1:item:a}\SGAInlineRefTarget{sga3:xii:5.1:item:a}(a)}] \emph{\(M\) 的每个在 \(S\) 上有限型的子预概形 \(U\)，都包含于一个在 \(S\) 上有限型的闭子预概形。\(M\) 的每个在 \(S\) 上有限型的闭子预概形 \(X\)，都在 \(S\) 上仿射。}
\item[\textup{\SGAInlineRefTarget{sga3:xii:5.1:item:b}(b)}] \emph{设 \(S\) 拟紧，\(X\) 是 \(M\) 的在 \(S\) 上有限型的闭子预概形。则存在整数 \(n>0\)，使对 \(n\) 的每个倍数 \(m\)，诱导态射
\[
  u_m|_X:X\longrightarrow T_m
\]
都是闭浸入。}
\end{enumerate}
\emph{证明。}
\textup{(a)} 对第一个断言，取 \(X\) 为 \(U\) 的概形闭包\textup{（EGA~I，9.5.1、9.5.3）}。这个闭包有定义，因为浸入 \(i:U\to M\) 拟紧：\(U\) 在 \(S\) 上有限型，而由 \SGARefLink{sga3:xi:theorem:4:1}{XI~4.1}，\(M\) 在 \(S\) 上分离。只须再证这样的 \(X\) 在 \(S\) 上仿射，这也同时证明第二个断言。问题在 \(S\) 上是局部的，所以可设 \(S\) 仿射。由于 \(U\) 在 \(S\) 上有限型，它包含于 \(M\) 的某个拟紧开集中。因此可归约到 \(U\) 本身是一个在 \(S\) 上有限型、从而拟紧的开子预概形的情形。注意，在 \(S\) 上仿射的概形的闭子概形仍在 \(S\) 上仿射。
只须证明这样的 \(U\) 包含于 \(M\) 的某个闭仿射子预概形。在此形式下，通常的归约过程立即将我们归约到 \(S\) 是 Noether 的情形。同一论证也适用于\textup{\SGARefLink{sga3:xii:5.1:item:b}{(b)}}，将其归约到 \(S\) 仿射且 Noether 的情形。
回顾 \SGARefLink{sga3:xi:theorem:4:1}{XI~4.1} 证明中所用的各 \(T_n\) 的逆极限 \(T\)。它是 \(S\) 上的仿射预概形，虽然不是有限型，但其局部环是 Noether 的，IX~3.7 证明开头已指出这一点。更明确地，若 \(t=(t_n)\in T\)，则转移态射 \(T_m\to T_n\) 平滑，且各 \(T_m\) 的维数有界。因此存在 \(n_0\)，使对 \(n_0\) 的每个倍数 \(n\)，\(T_n\to T_{n_0}\) 在 \(t_n\) 处平展。刚才引用的证明或 \SGARefLink{sga3:xi:corollary:3:11}{XI~3.11} 表明，典范态射
\[
  u:M\longrightarrow T
\]
是浸入，且在局部环上诱导同构。它不一定是开浸入，也不一定是拟紧态射。
设 \(U\) 是 \(M\) 的拟紧开集。我们要证明，它在 \(T\) 中的闭包包含于 \(M\)。若 \(X\) 是 \(U\) 在 \(M\) 中的概形闭包，这将使
\[
  u|_X:X\longrightarrow T
\]
成为闭浸入。于是 \(X\) 仿射，证明了\textup{\SGARefLink{sga3:xii:5.1:item:a}{(a)}}。由于 \(U\) 是 Noether 的，它只有有限多个不可约分支。每个处于 \(U\) 的闭包中的 \(t\in T\)，都是某点 \(x\in U\) 的特殊化。\(t\) 在 \(T\) 中的局部环是 Noether 的，故 \(t\) 在赋予诱导约化结构的 \(\overline{\{x\}}\) 中的局部环 \(A\) 也是 Noether 的。
从 Noether 整局部概形
\[
  S'=\operatorname{Spec}(A)
\]
到 \(T\) 的典范态射将其闭点映为 \(t\)，一般点映为 \(x\in M\)。必须证明 \(t\in M\)。必要时，用某个适当的完备局部 \(A\)-代数（在 \(A\) 上平坦）对一个极小素理想的商替换 \(A\)\textup{（EGA 0\textsubscript{III}，10.3.1）}，可设 \(A\) 完备且剩余域代数闭。由 EGA~II，7.1.7，甚至还可进一步归约到离散赋值环。将 \(S'\) 重新记为 \(S\)，便归约到以下引理。
\paragraph{\textbf{引理 5.2。}}
\SGARefTarget{sga3:xii:5.2}
\emph{设 \(S\) 是剩余域代数闭的完备 Noether 整局部环的谱。设 \(G\) 是平滑仿射 \(S\)-群预概形，\(\bigl(H(n)\bigr)_{n>0}\) 是 \(G\) 中以正整数为指标的乘法型子群系，\(\eta\) 是 \(S\) 的一般点。假设：}
\begin{enumerate}
\item[\textup{(a)}] \emph{若 \(m\) 是 \(n\) 的倍数，则
\[
  H(n)={}_nH(m);
\]
}
\item[\textup{\SGAInlineRefTarget{sga3:xii:5.2:item:b}(b)}] \emph{在 \(G_\eta\) 中存在乘法型子群 \(H_\eta\)，使得
\[
  H(n)_\eta={}_n(H_\eta)
\]
对每个 \(n>0\) 成立。}
\end{enumerate}
\emph{则在 \(G\) 中存在乘法型子群 \(H\)，使对每个 \(n>0\) 有 \(H(n)={}_nH\)。}
\emph{证明。}
对每个 \(n\)，令
\[
  C_n=\operatorname{Centr}_G\bigl(H(n)\bigr).
\]
由 \SGARefLink{sga3:xi:corollary:5:3}{XI~5.3}，这是 \(G\) 的闭子群预概形，在 \(S\) 上平滑。将正整数按整除关系排序，则各 \(C_n\) 构成 \(G\) 的闭子预概形的递减族。由于 \(G\) 是 Noether 的，此族定常。令 \(C\) 为充分可整除的 \(n\) 所对应的 \(C_n\) 的共同值。则 \(C\) 满足与 \(G\) 相同的条件，且所有 \(H(n)\) 实际上都是 \(C\) 的子群。因此，用 \(C\) 替换 \(G\)，可设所有 \(H(n)\) 都是 \(G\) 的中心子群。
令 \(s\) 为 \(S\) 的闭点，\(Z_s\) 为 \SGARefLink{sga3:xii:4.4}{4.4} 所定义的 \(G_s\) 的约化中心。由 \SGARefLink{sga3:xii:4.1}{4.1} 的定义，\(Z_s\) 包含每个 \(H(n)_s\)。由于 \(A\) 完备，由 \SGARefLink{sga3:xi:theorem:5:8}{XI~5.8}，\(Z_s\) 提升为 \(G\) 中的乘法型子群预概形 \(Z\)。断言 \(Z\) 包含每个 \(H(n)\)。事实上，\(H(n)\) 是中心的，因而在 \(G\) 中与 \(Z\) 交换。由此得到的同态
\[
  Z\times H(n)\longrightarrow G
\]
由 \SGARefLink{sga3:ix:theorem:6:8}{IX~6.8} 有像 \(K\)，它是 \(G\) 中包含 \(Z\) 的乘法型子群。只须再证 \(K=Z\)；这由 \(K_s=Z_s\) 及 \SGARefLink{sga3:ix:theorem:5:1}{IX~5.1}~bis 推出。
于是归约到将引理中的 \(G\) 换为 \(S\) 上有限型乘法型群 \(Z\) 的类似命题，不必要求 \(Z\) 在 \(S\) 上平滑。由于 \(A\) 完备且剩余域代数闭，由 \SGARefLink{sga3:x:corollary:3:3}{X~3.3} 和 \SGARefLink{sga3:xii:1.4}{1.4}，\(Z\) 可对角化，故形如 \(D_S(M)\)，其中 \(M\) 是有限生成交换群。因此，由 \SGARefLink{sga3:ix:proposition:2:11:item:i:01}{IX~2.11(i)}，\(Z\) 的每个乘法型子群 \(H\) 都可对角化，故形如 \(D_S(N)\)，其中由 \SGARefLink{sga3:viii:corollary:3:2}{VIII~3.2}，\(N\) 是 \(M\) 的商群。于是各 \(H(n)\) 对应于 \(M\) 的商群 \(M(n)\)，而存在子群 \(H\) 使对每个 \(n\) 有 \(H(n)={}_nH\)，等价于存在 \(M\) 的商群 \(N\)，使
\[
  M(n)=N/nN
\]
对每个 \(n\) 成立。这立即来自 \(G_\eta\) 的子群 \(H_\eta\) 的存在性，其中对每个 \(n\) 有 \(H(n)_\eta={}_n(H_\eta)\)。这证明了\SGARefLink{sga3:xii:5.2}{引理 5.2}，从而证明了\SGARefLink{sga3:xii:theorem:5-1:item:a}{定理 5.1(a)}。
\textup{\SGARefLink{sga3:xii:5.2:item:b}{(b)}} 已知
\[
  u|_X:X\longrightarrow T
\]
是闭浸入。容易得到，对所有充分大的 \(m\)，复合
\[
  u_m|_X:X\longrightarrow T_m
\]
也是闭浸入。下文不用这一事实，故略去证明细节。
\paragraph{\textbf{推论 5.3。}}
\SGARefTarget{sga3:xii:5.3}
\emph{沿用 \SGARefLink{sga3:xii:5.1}{5.1} 的记号，设 \(U\) 是 \(M\) 的既开又闭、在 \(S\) 上有限型的子预概形。则 \(U\) 赋予 \(M\) 诱导的结构后，在 \(S\) 上仿射。若 \(S\) 拟紧，则存在整数 \(n>0\)，使对 \(n\) 的每个倍数 \(m\)，诱导态射
\[
  u_m|_U:U\longrightarrow T_m
\]
既是开浸入又是闭浸入；等价地，它是到 \(T_m\) 的某个赋予诱导结构的既开又闭子预概形的同构。}
\emph{证明。}
第一个断言来自 \SGARefLink{sga3:xii:5.1:item:a}{5.1(a)}，第二个来自 \SGARefLink{sga3:xii:5.1:item:b}{5.1(b)}，因为由 \SGARefLink{sga3:xi:corollary:2:2}{XI~2.2}~bis，\(u_m:M\to T_m\) 平滑，而平滑浸入，即平展浸入，是开浸入。
\paragraph{\textbf{推论 5.4。}}
\SGARefTarget{sga3:xii:5.4}
\emph{设 \(S\) 是预概形，\(G\) 是约化秩局部常值的平滑仿射 \(S\)-群预概形\textup{（\SGARefLink{sga3:xii:1.7:item:b}{1.7(b)}）}，\(\mathcal T\) 是“\(G\) 的极大环面预概形”\textup{（\SGARefLink{sga3:xii:1.10}{1.10}）}。则 \(\mathcal T\) 在 \(S\) 上平滑且仿射。若 \(T\) 是 \(G\) 的极大环面，\(N(T)\) 为其正规化子，则 \(G/N(T)\)\textup{（\SGARefLink{sga3:xi:corollary:5:3}{XI~5.3}~bis）}在 \(S\) 上仿射。对 \(G/C(T)\) 也有同样结论，其中 \(C(T)\) 是 \(T\) 的中心化子，只要
\[
  W(T)=N(T)/C(T)
\]
在 \(S\) 上有限；参见\ \SGARefLink{sga3:xii:2.1:item:b}{2.1(b)}。}
\emph{证明。}
第二个断言包含在第一个断言中，因为由 \SGARefLink{sga3:xi:remark:5:5}{XI~5.5}~bis，共轭定理将 \(G/N(T)\) 与 \(\mathcal T\) 等同。对第一个断言，注意由构造，\(\mathcal T\) 同构于 \(M\) 的一个既开又闭子预概形。事实上，可设 \(G\) 的约化秩恒等于 \(r\)；则 \(\mathcal T\) 是 \(M\) 中对应于相对维数为 \(r\) 的子环面的子预概形，即 \(M\) 的最大子预概形，使其上万有乘法型子群
\[
  H\subset G_M
\]
是相对维数为 \(r\) 的环面。因此\SGARefLink{sga3:xii:5.3}{推论 5.3} 适用。最后，\(G/C(T)\) 在
\[
  G/N(T)\simeq \bigl(G/C(T)\bigr)/W(T),
\]
上有限，故因 \(G/N(T)\) 仿射而仿射。
\paragraph{\textbf{推论 5.5。}}
\SGARefTarget{sga3:xii:5.5}
\emph{设 \(G\) 是域 \(k\) 上的平滑仿射代数群。则 \(G\) 的乘法型子群概形 \(M\) 是仿射 \(k\)-概形的不交并。对 \(G\) 的每个乘法型子群 \(H\)，若 \(C\) 和 \(N\) 分别表示其中心化子和正规化子，则商 \(G/C\) 和 \(G/N\) 仿射。}
\emph{证明。}
由 \SGARefLink{sga3:xi:corollary:5:1}{XI~5.1}~bis，\(M\) 的每个有限闭子集在 \(G\) 于 \(M\) 的作用下的饱和集都是开的。事实上，可归约到 \(k\) 代数闭的情形，从而归约到某个 \(k\)-有理点 \(x\) 的轨道；此时，由所引结果，态射
\[
  G\longrightarrow M,\qquad g\longmapsto g\cdot x
\]
平滑，因而是开的，所以其像为开集。
令 \(U\) 为 \(M\) 的闭点在 \(G\) 作用下的等价类的并。则 \(U\) 开且包含 \(M\) 的每个闭点；由零点定理，\(U=M\)。因此 \(M\) 是开子预概形的不交并，这些子预概形必也是闭的。故 \(M\) 是各子预概形 \(M_i\) 的和预概形，每个 \(M_i\) 都是某个闭点的 \(G\)-轨道。每个 \(M_i\) 拟紧，因而有限型，故由 \SGARefLink{sga3:xii:5.3}{5.3} 仿射。
若 \(H\) 是 \(G\) 的乘法型子群，它对应于 \(M\) 的一个 \(k\)-有理点 \(x\)，且由 \SGARefLink{sga3:xi:remark:5:5}{XI~5.5}~bis，\(G/N\) 可与 \(x\) 在 \(G\) 下的轨道等同。因此 \(G/N\) 仿射。由于由 \SGARefLink{sga3:xi:proposition:5:9}{XI~5.9}，\(C\) 是 \(N\) 的开子群，商 \(G/C\) 在 \(G/N\) 上有限；事实上，
\[
  G/N\simeq (G/C)/(N/C).
\]
所以 \(G/C\) 也仿射。
此证明还给出如下结果。
\paragraph{\textbf{推论 5.6。}}
\SGARefTarget{sga3:xii:5.6}
\emph{在 \SGARefLink{sga3:xii:5.5}{5.5} 关于 \(G\) 和 \(H\) 的假设下，\(M\) 中由 \(G\) 的与 \(H\) 局部共轭的乘法型子群组成的子预概形 \(U\)，是 \(M\) 的既开又闭子预概形。}
形象地说，\(G\) 的每个作为与 \(H\) 共轭的子群之极限的乘法型子群 \(H'\)，本身也与 \(H\) 共轭。
\paragraph{\textbf{注 5.7。}}
\SGARefTarget{sga3:xii:5.7}
设 \(S\) 是预概形，\(G\) 是平滑仿射 \(S\)-群预概形，\(H\) 是有限型乘法型 \(S\)-预概形，置
\[
  M=\operatorname{Hom}_{S\text{-gr}}(H,G).
\]
由 \SGARefLink{sga3:xi:corollary:4:2}{XI~4.2}，\(M\) 可表，在 \(S\) 上平滑，并在 \(S\) 上分离。可以对 \(M\) 证明与 \SGARefLink{sga3:xii:5.1}{5.1} 完全类似的结果：或者通过类似于 \SGARefLink{sga3:xi:corollary:4:2}{XI~4.2} 的论证归约到 \SGARefLink{sga3:xii:5.1}{5.1}，或者仿照 \SGARefLink{sga3:xii:5.1}{5.1} 的证明直接论证。由此得到 \SGARefLink{sga3:xii:5.3}{5.3}、\SGARefLink{sga3:xii:5.5}{5.5} 和 \SGARefLink{sga3:xii:5.6}{5.6} 的相应变体；其陈述留给读者。
\section[极大环面与 Cartan 子群]
{不一定仿射的代数群的极大环面与 Cartan 子群\textup{（代数闭基域）}}
\SGARefTarget{sga3:xii:6}
\paragraph{\textbf{引理 6.1。}}
\SGARefTarget{sga3:xii:6.1}
\emph{设 \(G\) 是域 \(k\) 上的连通代数群，\(Z\) 是 \(G\) 的中心中有限指数的代数子群。则 \(G/Z\) 仿射。}
\emph{证明。}
可设 \(Z\) 是 \(G\) 的中心，因为仿射概形上的有限概形仿射。考虑向量空间
\[
  P^n=P^n(G)=
  \mathcal O_{G,e}/\mathfrak m_{G,e}^{\,n+1}
  \qquad (n\geq 0),
\]
其中 \(\mathfrak m_{G,e}\) 是单位元处的极大理想。群 \(G\) 通过伴随表示作用于各 \(P^n(G)\)。令 \(Z_n\) 为对应同态
\[
  \operatorname{ad}_n:G\longrightarrow \operatorname{GL}(P^n).
\]
的核。利用 \(G\) 的连通性，容易验证 \(Z\) 是各 \(Z_n\) 的交。由于 \(G\) 是 Noether 的，此交等于某个 \(Z_n\)。取商后，\(\operatorname{ad}_n\) 定义单态射
\[
  G/Z_n\longrightarrow \operatorname{GL}(P^n),
\]
它是闭浸入。因此 \(G/Z_n\) 仿射，\(G/Z\) 也仿射。
\paragraph{\textbf{引理 6.2。}}
\SGARefTarget{sga3:xii:6.2}
\emph{设 \(G\) 是域 \(k\) 上的平滑代数群，\(Z\) 是中心代数子群，置 \(G'=G/Z\)，令
\[
  u:G\longrightarrow G'
\]
为典范同态。设 \(T\) 是 \(G\) 的乘法型子群，\(T'=u(T)\)，\(C\) 是 \(T\) 在 \(G\) 中的中心化子，\(C'\) 是 \(T'\) 在 \(G'\) 中的中心化子。则
\[
  C'^{\,0}\subset u(C).
\]
}
\emph{证明。}
可设 \(k\) 代数闭。置
\[
  C_1=\bigl(u^{-1}(C'^{\,0})\bigr)_{\mathrm{red}}.
\]
只须证明 \(C_1\subset C\)。事实上，\(u(C_1)\) 连通且在 \(C'^{\,0}\) 中指数有限，故在集合意义上等于后者；而 \(C'\)，从而 \(C'^{\,0}\)，平滑，故在概形意义上也相等。
从 \(G\times G\) 到 \(G\) 的交换子态射
\[
  (g,t)\longmapsto gtg^{-1}t^{-1}
\]
诱导态射
\[
  \varphi:C_1\times T\longrightarrow Z_1,
  \qquad Z_1=Z_{\mathrm{red}}.
\]
事实上，定义域约化，且对 \(g\in C_1(k)\) 及 \(t\in T(k)\)，元素 \(gtg^{-1}t^{-1}\) 属于 \(Z(k)\)，因为 \(g\) 模 \(Z\) 中心化 \(T\)。
在 \(k\)-有理点上计算可见，\(\varphi(g,t)\) 关于 \(g\) 和 \(t\) 都是可加的，因而是“双线性”的。我们断言，此同态恒为零；这就证明 \(C_1\subset C\)。由 \(T\) 的稠密性定理，归约到 \(T\) 有限的情形，于是对某个整数 \(n>0\) 有 \(n\operatorname{id}_T=0\)。此时态射 \(\varphi\) 由群同态
\[
  C_1\longrightarrow \operatorname{Hom}_{k\text{-gr}}(T,Z_1).
\]
定义。由于 \(Z_1\) 交换且平滑，右边的群可由 \(k\) 上的平展代数群表示。由于 \(C_1\) 连通，从 \(C_1\) 到该群的每个群同态都为零。
\paragraph{\textbf{推论 6.3。}}
\SGARefTarget{sga3:xii:6.3}
\emph{在上述假设下，设 \(C'\) 连通。则
\[
  C'=u(C),
  \qquad
  C=u^{-1}(C').
\]
}
\emph{证明。}
此时 \(C'=C'^{\,0}\)。此外，\(C\) 显然包含 \(Z\)，故
\[
  C=u^{-1}\bigl(u(C)\bigr).
\]
\paragraph{\textbf{引理 6.4。}}
\SGARefTarget{sga3:xii:6.4}
\emph{设 \(G\) 是域 \(k\) 上的代数群，\(Z\) 是中心代数子群，使得
\[
  G/Z=G'
\]
是环面。则 \(G\) 交换。若 \(k\) 代数闭，则 \(G\) 中存在环面 \(T\)，使 \(u(T)=G'\)，其中
\[
  u:G\longrightarrow G/Z=G'
\]
为典范同态。}
\emph{证明。}
可设 \(k\) 代数闭。再次考虑交换子态射
\[
  G\times G\longrightarrow G,\qquad
  (g,h)\longmapsto ghg^{-1}h^{-1}.
\]
由于 \(Z\) 是中心的，此态射经由态射 \(G'\times G'\to G\) 分解。由于 \(G'\) 交换，后者取值于 \(Z\)，甚至取值于 \(Z_{\mathrm{red}}\)，因为 \(G'\times G'\) 约化。与前面一样，得到的态射
\[
  \varphi:G'\times G'\longrightarrow Z_{\mathrm{red}}
\]
双线性，故为零。因此 \(G\) 交换。
为在 \(G\) 中找到满足 \(u(T)=G'\) 的环面 \(T\)，可用 \(G_{\mathrm{red}}^{\,0}\) 替换 \(G\)，从而设 \(G\) 平滑且连通。用 \(Z_{\mathrm{red}}^{\,0}\) 替换 \(Z\)，也可设 \(Z\) 平滑且连通。由 Chevalley 的一个著名定理，\footnote{编者注：此处应补充参考文献。}\(Z\) 是阿贝尔簇按仿射群的扩张。因此归约到 \(Z\) 仿射或 \(Z\) 为阿贝尔簇这两种情形。第一种情形中 \(G\) 仿射，结果是经典的\textup{（BIBLE~7，定理 3(a)）}。
于是设 \(Z\) 是阿贝尔簇。从加法群 \(\mathbb G_a\) 到环面 \(G'=R\)，或到阿贝尔簇 \(Z\)，每个同态都平凡。因此，从 \(\mathbb G_a\) 到 \(G\) 的每个同态也平凡；所以 \(G\) 不含同构于 \(\mathbb G_a\) 的子群。Chevalley 结构定理给出正合列
\[
  0\longrightarrow T\longrightarrow G\longrightarrow A
  \longrightarrow 0, \tag{*}\label{sga3:xii:lemma:6-4:equation:star}
\]
其中 \(A\) 是阿贝尔簇，\(T\) 是平滑连通仿射群。由于 \(T\) 交换且不含加法子群，\(T\) 是环面；它显然是 \(G\) 的唯一极大环面。只须再证，每个满态射
\[
  u:G\longrightarrow R
\]
到环面 \(R\) 都满足 \(u(T)=R\)。
置
\[
  \operatorname{Hom}_{\mathrm{gr}}(T,\mathbb G_m)=M,
  \qquad
  \operatorname{Hom}_{\mathrm{gr}}(R,\mathbb G_m)=P.
\]
它们是有限生成自由 \(\mathbb Z\)-模，且
\[
  T=D_k(M),\qquad R=D_k(P).
\]
令
\[
  B=\operatorname{Ext}^{1}_{k\text{-gr}}(A,\mathbb G_m);
\]
这也是阿贝尔簇 \(A\) 的对偶阿贝尔簇的 \(k\)-有理点集。有
\[
  \operatorname{Ext}^{1}(A,T)
    =\operatorname{Hom}_{\mathrm{gr}}(M,B),
  \qquad
  \operatorname{Ext}^{1}(A,R)
    =\operatorname{Hom}_{\mathrm{gr}}(P,B). \tag{**}\label{sga3:xii:lemma:6-4:equation:double-star}
\]
特别地，\(A\) 按 \(T\) 的扩张 \(G\) 由一个同态
\[
  \theta:M\longrightarrow B.
\]
确定。
正合列\textup{(*)} 还给出
\[
  0\longrightarrow \operatorname{Hom}(A,R)
  \longrightarrow \operatorname{Hom}(G,R)
  \longrightarrow \operatorname{Hom}(T,R)
  \longrightarrow \operatorname{Ext}^{1}(A,R).
\]
此外，\(\operatorname{Hom}(A,R)=0\)，而最后一个映射可与由 \(\theta:M\to B\) 诱导的同态
\[
  \operatorname{Hom}(P,M)\longrightarrow
  \operatorname{Hom}(P,B)
\]
等同。置
\[
  N=\ker\theta,
\]
得到典范双射
\[
  \operatorname{Hom}(G,R)
  \simeq \operatorname{Hom}(P,N)
  \simeq \operatorname{Hom}(S,R),
  \qquad S=D_k(N).
\]
这立即有如下几何描述。
\paragraph{\textbf{引理 6.5。}}
\SGARefTarget{sga3:xii:6.5}
\emph{设 \(G\) 是代数闭域 \(k\) 上阿贝尔簇 \(A\) 按环面 \(T=D_k(M)\) 的扩张，由同态
\[
  \theta:M\longrightarrow
  B=\operatorname{Ext}^{1}(A,\mathbb G_m),
\]
定义。置 \(N=\ker\theta\)，令 \(S=D_k(N)\) 为对应的环面，它同构于 \(T/U\)，其中 \(U=D_k(M/N)\)。考虑 \(A\) 按 \(S\) 的扩张 \(H=G/U\)。此扩张分裂，\footnote{编者注：此扩张分裂。}故有从 \(H\) 到 \(S\) 的唯一投影，从而有延拓典范同态 \(T\to S\) 的唯一同态
\[
  v:G\longrightarrow S
\]
对每个环面 \(R\) 和每个同态 \(u:G\to R\)，存在唯一同态 \(u':S\to R\)，使 \(u=u'v\)。特别地，
\[
  \operatorname{Im}(u)=\operatorname{Im}(u')=u(T).
\]
}
特别可得，若 \(u\) 是满态射，则其在 \(T\) 上的限制也是满态射。\SGARefLink{sga3:xii:6.4}{6.4} 的证明至此完成。
\paragraph{\textbf{定理 6.6。}}
\SGARefTarget{sga3:xii:6.6}
\emph{设 \(G\) 是代数闭域 \(k\) 上的平滑连通代数群。}
\begin{enumerate}
\item[\textup{\SGAInlineRefTarget{sga3:xii:theorem:6-6:item:a}\SGAInlineRefTarget{sga3:xii:6.6:item:a}(a)}] \emph{\(G\) 的极大环面彼此共轭。}
\item[\textup{\SGAInlineRefTarget{sga3:xii:6.6:item:b}\SGAInlineRefTarget{sga3:xii:theorem:6-6:item:b}(b)}] \emph{若 \(T\) 是 \(G\) 的环面，则其中心化子平滑且连通。}
\item[\textup{\SGAInlineRefTarget{sga3:xii:6.6:item:c}\SGAInlineRefTarget{sga3:xii:theorem:6-6:item:c}(c)}] \emph{对应
\[
  T\longmapsto \operatorname{Centr}_G(T)
\]
给出 \(G\) 的极大环面 \(T\) 与第 1 节定义的 \(G\) 的 Cartan 子群之间的双射。\(G\) 的代数子群 \(C\) 是 Cartan 子群，当且仅当它平滑、幂零，且在集合意义上等于其连通正规化子；此时它甚至在概形意义上等于其连通正规化子。此外，
\[
  C=\operatorname{Centr}_G(T),
\]
其中 \(T\) 是 \(C\) 的唯一极大环面，且
\[
  \operatorname{Norm}_G(C)=\operatorname{Norm}_G(T).
\]
}
\item[\textup{\SGAInlineRefTarget{sga3:xii:6.6:item:d}(d)}] \emph{设 \(v:G\to H\) 是平滑连通代数群的满态射。\(H\) 的极大环面（分别为 Cartan 子群）是 \(G\) 的极大环面（分别为 Cartan 子群）的像。若 \(T\) 是 \(G\) 的极大环面，\(C\) 是其中心化子，则 \(v(C)\) 是 \(v(T)\) 的中心化子。}
\item[\textup{\SGAInlineRefTarget{sga3:xii:theorem:6-6:item:e}\SGAInlineRefTarget{sga3:xii:6.6:item:e}(e)}] \emph{在\textup{\SGARefLink{sga3:xii:corollary:6-7:equation:d}{(d)}} 的假设下，设 \(\ker v\) 是 \(G\) 的中心子群。则
\[
  T\longmapsto v(T),
  \qquad
  C\longmapsto v(C)
\]
给出 \(G\) 的极大环面（分别为 Cartan 子群）与 \(H\) 的极大环面（分别为 Cartan 子群）之间的双射。\(G\) 的 Cartan 子群包含 \(G\) 的中心，因而包含 \(\ker v\)；它们恰好是形如 \(v^{-1}(C')\) 的群，其中 \(C'\) 是 \(H\) 的 Cartan 子群。}
\end{enumerate}
还记下\SGARefLink{sga3:xii:theorem:6-6:item:c}{定理 6.6(c)} 的如下直接推论。
\paragraph{\textbf{推论 6.7。}}
\SGARefTarget{sga3:xii:6.7}
\emph{群 \(G\) 幂零，即 \(G(k)\) 幂零，当且仅当 \(G\) 中每个环面都是中心的；又等价于 \(G\) 存在唯一极大环面。此时，此环面是 \(G\) 的最大子环面。}
\begin{proof}[\SGARefLink{sga3:xii:6.6}{定理 6.6} 的证明]
设 \(Z\) 是 \(G\) 的中心代数子群，置 \(G'=G/Z\)，令
\[
  u:G\longrightarrow G'
\]
为典范同态。则 \(G'\) 平滑且连通。若 \(T'\) 是 \(G'\) 的子环面，由 \SGARefLink{sga3:xii:6.4}{6.4}，\(u^{-1}(T')\) 交换，且 \(T'\) 是 \(u^{-1}(T')\) 的某个子环面的像，从而是 \(G\) 的某个子环面的像。由于 \(u^{-1}(T')\) 交换，它显然有一个最大子环面 \(T\)：两个子环面之和是包含两者的第三个子环面。因此
\[
  u(T)=T'.
\]
立即可得，对 \(G\) 的每个极大环面 \(T\)，其像 \(T'=u(T)\) 是 \(G'\) 的极大环面，且
\[
  T\longmapsto u(T)
\]
是 \(G\) 与 \(G'\) 的极大环面之间的双射。
现在取
\[
  Z=\operatorname{Centr}(G)_{\mathrm{red}}^{0}.
\]
则由 \SGARefLink{sga3:xii:6.1}{6.1}，\(G'\) 仿射。其极大环面彼此共轭，故 \(G\) 的极大环面也彼此共轭，这证明了\textup{\SGARefLink{sga3:xii:6.6:item:a}{(a)}}。此外，\(G\) 幂零，或存在唯一极大环面，当且仅当 \(G'\) 分别具有相同性质。由于 \(G'\) 仿射，关于 \(G'\) 的这两个条件等价\((\mathrm{BIBLE}\ 6,\ \text{定理}\ 4,\ \text{推论}\ 2)\)；故对 \(G\) 也等价。若 \(G\) 存在唯一极大环面 \(T\)，则 \(T\) 不变，因而是中心的；又由于 \(G\) 的每个环面都包含于极大环面，每个环面都为中心的。反之，若每个环面都是中心的，极大环面也如此，共轭定理\textup{\SGARefLink{sga3:xii:6.6:item:a}{(a)}} 表明它只能有一个。这证明了 \SGARefLink{sga3:xii:6.7}{6.7}。
设 \(T\) 是 \(G\) 的任意环面，\(T'=u(T)\)，置
\[
  C'=\operatorname{Centr}_{G'}(T').
\]
群 \(C'\) 连通\((\mathrm{BIBLE}\ 6,\ \text{定理}\ 6\textup{(a)})\)。由 \SGARefLink{sga3:xii:6.3}{6.3}，\(T\) 的中心化子 \(C\) 是 \(u^{-1}(C')\)，因而由于 \(Z\) 连通，\(C\) 连通。这证明了\textup{\SGARefLink{sga3:xii:6.6:item:b}{(b)}}。若 \(T\) 极大，则 \(T'\) 极大；故 \(C'\) 幂零，而 \(C\) 是 \(C'\) 的中心扩张，因而也幂零。此外，\(T\) 是 \(C\) 的极大环面，由 \SGARefLink{sga3:xii:6.7}{6.7}，它是 \(C\) 的唯一极大环面。因此
\[
  T\longmapsto\operatorname{Centr}_G(T)
\]
是从 \(G\) 的极大环面到 \(G\) 的 Cartan 子群的双射。还有
\[
  \operatorname{Centr}_G(T)=C
  \subset \operatorname{Norm}_G(C)
  \subset \operatorname{Norm}_G(T).
\]
\(T\) 的中心化子在其正规化子中指数有限\((\text{参见 XI~5.9})\)；该处的论证在没有仿射性假设时仍有效，因为只用了所涉及的两个函子的可表性，\SGARefLink{sga3:xi:unit:6:5}{XI~6.5} 已指出这一点。由于 \(C\) 平滑且连通，可得
\[
  C=\operatorname{Norm}_G(C)^{0}.
\]
此外，极大环面与 Cartan 子群之间的双射表明，\(C\) 与 \(T\) 的正规化子具有相同的 \(k\)-值点。由于后一个正规化子平滑，
\[
  \operatorname{Norm}_G(T)=\operatorname{Norm}_G(C).
\]
为完成\textup{\SGARefLink{sga3:xii:6.6:item:c}{(c)}} 的证明，只须再证：\(G\) 的平滑连通幂零子群 \(C\)，若在其正规化子中指数有限，则是 Cartan 子群。由于 \(Z\) 是中心的，\(C\) 的正规化子包含 \(Z\)；而由于 \(Z\) 平滑且连通，有 \(Z\subset C\)。故
\[
  C=u^{-1}(C'),\qquad C'=u(C),
\]
且
\[
  \operatorname{Norm}_G(C)
  =u^{-1}\bigl(\operatorname{Norm}_{G'}(C')\bigr).
\]
因此 \(C'\) 连通、幂零，且在其正规化子中指数有限。由 \(\mathrm{BIBLE}\ 7\) 的定理 1，它是 \(G'\) 的 Cartan 子群，立即可知 \(C\) 是 \(G\) 的 Cartan 子群。
接着证明\textup{\SGARefLink{sga3:xii:6.6:item:e}{(e)}}。我们处于证明开头考虑的情形，其中 \(\ker u=Z\)、\(G'=H\)，且已知 \(T\mapsto u(T)\) 是 \(G\) 与 \(G'\) 的极大环面之间的双射。利用刚才证明的极大环面与 Cartan 子群之间的双射，得到 \(G\) 的 Cartan 子群 \(C\) 与 \(G'\) 的 Cartan 子群 \(C'\) 之间的双射：对于
\[
  C=\operatorname{Centr}_G(T)
\]
联系
\[
  C'=\operatorname{Centr}_{G'}(T'),
  \qquad T'=u(T).
\]
由于\textup{\SGARefLink{sga3:xii:6.6:item:b}{(b)}}，\(C'\) 连通，\SGARefLink{sga3:xii:6.3}{6.3} 给出
\[
  C'=u(C),\qquad C=u^{-1}(C'),
\]
这就证明了\textup{\SGARefLink{sga3:xii:6.6:item:e}{(e)}}。
只须再证\textup{\SGARefLink{sga3:xii:corollary:6-7:equation:d}{(d)}}。设 \(T\) 是 \(G\) 的极大环面。先证明 \(v(T)\) 是 \(H\) 的极大环面；由共轭定理\textup{\SGARefLink{sga3:xii:6.6:item:a}{(a)}}，这将蕴含 \(H\) 的全部极大环面都由此得到。设 \(R\) 是 \(H\) 中包含 \(v(T)\) 的环面。用 \(R\) 替换 \(H\)，用 \(v^{-1}(R)_{\mathrm{red}}^{0}\) 替换 \(G\)，可设 \(R=H\)，从而 \(H\) 是环面；只须再证 \(v(T)=H\)。再次置
\[
  Z=\operatorname{Centr}(G)_{\mathrm{red}}^{0},
  \qquad G'=G/Z,\qquad H'=H/v(Z).
\]
有各行正合的交换图
\begin{equation}
\tag{D}\label{sga3:xii:corollary:6-7:equation:d}
\begin{tikzcd}[column sep=large,row sep=large]
e \arrow[r] &
Z \arrow[r] \arrow[d,"v''"'] &
G \arrow[r,"u"] \arrow[d,"v"'] &
G' \arrow[r] \arrow[d,"v'"'] &
e \\
e \arrow[r] &
v(Z) \arrow[r] &
H \arrow[r,"u'"'] &
H' \arrow[r] &
e .
\end{tikzcd}
\end{equation}
已知 \(u(T)\) 是 \(G'\) 的极大环面。由于 \(G'\) 仿射，满态射
\[
  v':G'\longrightarrow H'
\]
是平滑连通仿射群之间的满态射。因此，\(v'(u(T))\) 是 \(H'\) 的极大环面\((\mathrm{BIBLE}\ 7,\ \text{定理}\ 3\textup{(a)})\)，故等于 \(H'\)。于是 \(u'(v(T))=H'\)，所以为证明 \(v(T)=H\)，只须证明 \(v(T)\supset v(Z)\)。现在，\(v(Z)\) 是 \(H\) 的平滑连通子群，因而是环面，且
\[
  v'':Z\longrightarrow v(Z)
\]
是满态射。由 \SGARefLink{sga3:xii:6.4}{6.4}，
\[
  v(Z)=v''(S)
\]
其中 \(S\) 是 \(Z\) 的环面。由于 \(S\) 在 \(G\) 中是中心的，它包含于极大环面 \(T\)，因此 \(v(Z)\subset v(T)\)。这证明了极大环面的\textup{\SGARefLink{sga3:xii:6.6:item:d}{(d)}}。
利用极大环面与 Cartan 子群之间的双射，只须再证：若 \(T\) 是 \(G\) 的极大环面，\(C\) 是其中心化子，则 \(v(C)\) 是 \(v(T)\) 的中心化子。回到图\textup{(D)}，现在不再假设 \(H\) 为环面，并置
\[
  T'=u(T),\qquad C'=u(C).
\]
由\textup{\SGARefLink{sga3:xii:6.6:item:e}{(e)}}，\(C'\) 是极大环面 \(T'\) 的中心化子。由于 \(G'\)，从而 \(H'\)，仿射，\(v'(C')\) 是 \(H'\) 的极大环面 \(v'(T')\) 的中心化子\((\mathrm{BIBLE}\ 7,\ \text{定理}\ 3\textup{(a)})\)。等价地，\(u'(v(C))\) 是 \(u'(v(T))\) 的中心化子。由于 \(C\) 包含 \(Z\)，群 \(v(C)\) 包含 \(v(Z)\)。因此，它是 \(u'(v(C))\)，即 \(u'(v(T))\) 的中心化子在 \(u'\) 下的逆像。将\textup{\SGARefLink{sga3:xii:6.6:item:e}{(e)}} 用于 \(u':H\to H'\) 及 \(H\) 的极大环面 \(v(T)\)，可知此逆像恰好是 \(v(T)\) 的中心化子。\SGARefLink{sga3:xii:6.6}{6.6} 的证明至此完成。
\end{proof}
